\documentclass[11pt,a4paper]{article}

\usepackage[margin=1in]{geometry}
\usepackage{parskip}
\usepackage{titlesec}
\usepackage{enumitem}
\usepackage{hyperref}

\titleformat{\section}{\large\bfseries}{\thesection.}{0.5em}{}
\titleformat{\subsection}{\normalsize\bfseries}{\thesubsection}{0.5em}{}
\setlist{itemsep=2pt, topsep=4pt}

\hypersetup{colorlinks=true, linkcolor=black, urlcolor=blue}

\newcommand{\doctitle}[2]{
  \begin{center}
    {\LARGE\bfseries #1}\\[4pt]
    {\large #2}\\[6pt]
    Sahithyan Kandathasan \\
    Index No: 230557T
  \end{center}
  \vspace{1em}
}


\begin{document}

\doctitle{Monopoly Market Failure}{Case Study: The De Beers Diamond Cartel}

\section{Introduction}
A monopoly (or a cartel acting as one) can restrict competition and give a firm control over the price and quantity of a good far beyond what a competitive market would allow. When this power is used to deliberately limit supply, it can cause market failure by preventing prices from reflecting real production costs and consumer demand.

This case study examines De Beers Consolidated Mines and the global rough-diamond market it controlled for most of the twentieth century. For much of the past century, De Beers held a dominant position, at times controlling an estimated 80--90\% of the world's rough diamond trade, and used this position to deliberately restrict supply and manage prices.

\section{Background of De Beers}
De Beers was founded in South Africa in 1888 and grew, through acquisitions and agreements with other diamond-producing regions, into the central controlling force of the global diamond trade. It did not simply mine diamonds; it also acted as the main channel through which diamonds from many other producers (including mines in Russia, Central Africa, and elsewhere) reached the market, via its Central Selling Organisation.

Diamonds are a naturally scarce resource, but not scarce enough on their own to justify the extremely high prices maintained for most of the twentieth century. De Beers' dominance came from controlling the supply chain itself, not merely from owning the largest mines.

\section{Why Was De Beers Dominant?}
De Beers built its dominance by acquiring or forming agreements with diamond producers across multiple countries, so that most of the world's rough diamonds passed through its Central Selling Organisation before reaching cutters and retailers. This let it act as a single chokepoint for global diamond supply, even where it did not own the mines directly.

As a result, De Beers could decide how many diamonds reached the market in a given period, largely independent of how many were actually mined. Having this much control over supply is not automatically illegal everywhere, but the concern arises from what a firm does with that control.

\section{The Supply Restriction Problem}
De Beers used its position to manage the global diamond supply deliberately, including:
\begin{itemize}
  \item Stockpiling large quantities of diamonds rather than releasing them to market.
  \item Buying up diamonds from new discoveries or competing producers to prevent independent sales.
  \item Coordinating with producing countries and other firms to limit total output.
  \item Controlling the pace at which diamonds were released to cutters and wholesalers.
\end{itemize}

This was a problem because it meant the quantity of diamonds available to the market was set by De Beers' commercial strategy rather than by genuine scarcity or production capacity, keeping prices well above what a competitive market would produce.

\section{How Did This Cause Market Failure?}
The market failure can be explained as follows:

\begin{center}
Dominant control of the supply chain\\
$\downarrow$\\
Deliberate restriction and stockpiling of diamonds\\
$\downarrow$\\
Artificial scarcity, independent of actual mining output\\
$\downarrow$\\
Prices held far above competitive levels\\
$\downarrow$\\
Reduced consumer surplus and blocked independent producers\\
$\downarrow$\\
Inefficient allocation of resources
\end{center}

In a competitive market, an increase in supply from new mine discoveries would normally push prices down, and independent producers could sell freely at competitive prices. Here, De Beers absorbed new supply into its own stockpiles and controlled release, so prices did not fall to reflect the true abundance of diamonds. The result was an inefficient market outcome: consumers paid more than a competitive price would justify, and independent producers who tried to sell outside the De Beers channel faced a firm with the power to flood the market and undercut them.

\section{Impact on the Market}

\textbf{Impact on Consumers}

Consumers paid significantly inflated prices for diamonds relative to their actual scarcity and production cost, sustained in part by coordinated marketing (such as the "A Diamond is Forever" campaign) that reinforced demand alongside the restricted supply.

\textbf{Impact on Competition}

Independent and newly discovered diamond sources faced strong pressure to sell through De Beers' channel rather than compete directly, since De Beers could use its stockpiles to depress prices for any producer who sold independently.

\textbf{Impact on Producing Countries}

Some diamond-producing countries and smaller mining operations had limited bargaining power over pricing and output decisions that were effectively made by De Beers rather than by individual market participants.

\textbf{Impact on Market Efficiency}

With supply deliberately decoupled from actual production and scarcity, prices failed to act as an accurate signal of the diamond market's real conditions, leading to a persistently inefficient allocation of resources.

\section{Solutions I Would Propose}
If I were the decision maker at the time, I would propose the following solutions:

\begin{enumerate}
  \item \textbf{Antitrust Enforcement on Supply Coordination}\\
  Apply competition law specifically to stockpiling and output-restriction agreements between diamond producers, not only to price-fixing agreements, since supply control achieves the same effect.

  \item \textbf{Support Independent Sales Channels}\\
  Help independent and newly discovered producers access buyers directly (for example through open auctions), reducing their dependence on a single controlling intermediary.

  \item \textbf{Transparency in Supply and Stockpiles}\\
  Require reporting of production volumes and stockpile levels so that regulators and the market can assess whether supply is being deliberately withheld.

  \item \textbf{Encourage Alternative Supply Sources}\\
  Support the development of diamond-producing regions and channels outside the dominant firm's control, increasing the number of independent competitors over time.

  \item \textbf{Consumer Awareness}\\
  Support independent information about actual diamond scarcity and pricing, to counter marketing narratives that reinforced artificial scarcity.
\end{enumerate}

\section{Later Developments}
Over subsequent decades, De Beers' control gradually weakened as new producers (including Russian and Canadian mines) increasingly sold outside its channel, as competition regulators in various jurisdictions restricted its practices, and as the rise of lab-grown diamonds introduced a genuine substitute product. De Beers itself scaled back its stockpiling model and shifted toward a smaller market share operating under closer regulatory scrutiny.

\section{Conclusion}
The De Beers case demonstrates how a dominant firm's control over supply, rather than price-fixing alone, can cause market failure. By restricting the flow of a naturally abundant resource to the market, De Beers sustained prices far above competitive levels for decades. The case shows the importance of monitoring supply-side control, not just pricing behaviour, and of supporting independent producers and transparent markets to prevent this kind of dominance from persisting.

\section*{References}
\begin{enumerate}
  \item Epstein, E.J. (1982). \textit{The Rise and Fall of Diamonds: The Shattering of a Brilliant Illusion}.
  \item Federal Trade Commission -- historical antitrust actions concerning De Beers and the U.S. diamond trade.
  \item European Commission -- competition proceedings concerning De Beers' rough diamond supply agreements (2000s).
\end{enumerate}

\end{document}
